\scanpage{378}
\mspage{\deleted{\unsure{116}}117}
% archive sheet number: 377
associant à \deleted{\unsure{tout|\ensuremath{\mapsto}}} $a\in X$ le foncteur \uline{\unsure{const.}} \inserted{sur $Y$} de valeur $a$, \unsure{en|on} \deleted{\unsure{prenant l'image inverse du}} \inserted{\unsure{prenant|posant}} \deleted{\unsure{groupoïde}}
\[
H_T(a,b)\simeq\mathrm{Hom}_{X(Y)\Sigma_{\unsure{Y|y}}^{-1}}(a_Y,b_Y).
\]

Dans le $=$ cas géométrique, je me suis convaincu d'ailleurs que $C_\infty$ était, en tant qu'objet de $\mathrm{Hot}$ (pas seulement de $\mathrm{Hot}_0^{\unsure{\wedge|1}}$) \uline{unique à \inserted{$\mathrm{Hot}$-}isomorphisme près}, mais (chose étrange) peut-être pas à $\mathrm{Hot}$-isomorphisme \uline{unique} près.

De façon générale, revenant aux conditions du début du \S, $C_\infty$ doit s'écrire $C_\infty(\underline t)$, elle dépend du choix du foncteur cofinal
\[
\underline t:\mathbb N^\circ\longrightarrow T,
\]
Quand $\underline t$ varie dans la catégorie
\[
\underline{\mathrm{Hom\,cof}}(\mathbb N^\circ,T),
\]
les $\underline C_\infty(\underline t)$ varient de façon contravariante en $\underline t$ (donc définissent une catégorie fibrée sur $\underline{\mathrm{Hom}}\ \underline{\mathrm{cof}}(\mathbb N^\circ,T)$). De plus, le foncteur
\[
C_\infty(\underline t')\mathrel{\unsure{-|\to}}\underline C_\infty(\underline t)
\]
associé à un morphisme de systèmes projectifs
\[
\underline t\mathrel{-}\underline t'
\]
est un \uuline{iso.} dans \uline{$\mathrm{Hot}$}. Donc \deleted{les} le
